\subsection{Формирование функциональных и нефункциональных требований}

Требования к проектируемой системе определяются сценариями использования
ончейн-данных прикладными сервисами и особенностями самой предметной области.
В частности, система должна обеспечивать получение данных от доверенного
источника, воспроизводимую обработку подтвержденной цепи, корректную работу с
неподтвержденными транзакциями, устойчивость к reorg и удобное предоставление
результатов обработки внешним системам.

Для дальнейшего проектирования требования целесообразно разделить на
функциональные, требования к согласованности данных, эксплуатационные,
требования к безопасности и требования к расширяемости. Такая классификация
позволяет не только описать ожидаемые возможности системы, но и зафиксировать
критерии, которым должны удовлетворять архитектурные и программные решения.

\begin{table}
    \caption{Основные группы требований к системе}
    \begin{tabular}{|p{0.34\textwidth}|p{0.54\textwidth}|}\hline
    Группа требований & Содержание \\ \hline
    Функциональные & Индексация блоков, транзакций, входов, выходов, UTXO, адресных балансов, mempool-данных, управление заданиями индексирования, выдача данных через API \\ \hline
    Требования к согласованности & Идемпотентность записи, транзакционность операций, корректная обработка reorg, детерминированное обновление производных агрегатов \\ \hline
    Эксплуатационные & Контейнерное развертывание, логирование, публикация метрик, восстановление после перезапуска без потери критического состояния \\ \hline
    Требования к безопасности & Аутентификация административных операций, защищенное подключение к RPC, управление секретами через внешнюю конфигурацию и переменные окружения \\ \hline
    Требования к расширяемости & Возможность добавления новых модулей и развития системы без нарушения базовых контрактов взаимодействия \\ \hline
    \end{tabular}
    \label{tab:requirements-groups}
\end{table}

Функциональные и нефункциональные требования образуют основу всех дальнейших
проектных решений. В частности, требование stateless-характера исполнения ведет к
вынесению критического состояния в PostgreSQL, а требование воспроизводимого
развертывания обосновывает использование Docker и Docker Compose
\cite{docker_docs, compose_docs}. Требования к прикладной выдаче данных
обусловливают необходимость построения специализированной модели хранения,
ориентированной не только на запись блокчейн-событий, но и на быстрые выборки
по адресам, транзакциям, блокам и заданиям индексирования.
